% together/microopt.tex
% mainfile: ../perfbook.tex
% SPDX-License-Identifier: CC-BY-SA-3.0

\section{微优化}
\label{sec:together:Micro-Optimization}
%
\epigraph{魔鬼藏在细节中。}{佚名}

本书中展示的数据结构都是直接编码的，
没有针对底层系统的缓存层次结构进行调整。
此外，许多实现使用函数指针来进行
键到哈希的转换和其他频繁操作。
虽然这种方法提供了简洁性和可移植性，
但在许多情况下确实会损失一些性能。

以下各节涉及特化、内存节约和硬件考虑。
请不要把这些简短的章节误认为
是对此主题的权威论述。
整本书都可以写关于针对特定 CPU 的优化，
更不用说当今常用的各种 CPU 系列了。

\subsection{特化}
\label{sec:together:Specialization}

\cref{sec:datastruct:Non-Partitionable Data Structures}
（\cpageref{sec:datastruct:Non-Partitionable Data Structures}）
中介绍的可调整大小的 hash table 使用了
不透明类型作为键。
这允许极大的灵活性，允许使用任何类型的键，
但由于通过函数指针调用也会产生显著的开销。
现代硬件使用复杂的分支预测技术来最小化此开销，
但另一方面，现实世界的软件通常比
今天大型硬件分支预测表所能容纳的更大。
对于通过指针的调用尤其如此，
在这种情况下，分支预测硬件除了
分支是否跳转的信息外，还必须记录一个指针。

可以通过将 hash table 实现特化为给定的
键类型和哈希函数来消除此开销，
例如使用 C++ 模板。
这样做可以消除
\cref{lst:datastruct:Resizable Hash-Table Data Structures}
（\cpageref{lst:datastruct:Resizable Hash-Table Data Structures}）
中 \co{ht} 结构中的 \co{->ht_cmp()}、\co{->ht_gethash()} 和
\co{->ht_getkey()} 函数指针。
它还消除了通过这些指针的相应调用，
这可以允许编译器内联生成的固定函数，
不仅消除了调用指令的开销，
还消除了参数编组的开销。

\QuickQuiz{
	这些特化实际上能节省多少？
	它们真的值得吗？
}\QuickQuizAnswer{
	第一个问题的答案留给读者作为练习。
	尝试特化可调整大小的 hash table，
	看看能带来多少性能提升。
	第二个问题无法笼统地回答，
	而必须针对特定的用例来回答。
	某些用例对性能和可扩展性极为敏感，
	而其他用例则不那么敏感。
}\QuickQuizEnd

抛开这些不谈，与我在 1970 年代初刚开始
学习编程时可用的硬件相比，
现代硬件的一大优势是所需的特化要少得多。
这比在四千字节地址空间时代
允许了更高的生产力。

\subsection{位与字节}
\label{sec:together:Bits and Bytes}

本章讨论的 hash table 几乎没有尝试节省内存。
例如，
\cref{lst:datastruct:Resizable Hash-Table Data Structures}
（\cpageref{lst:datastruct:Resizable Hash-Table Data Structures}）
中 \co{ht} 结构的 \co{->ht_idx} 字段
始终只有零或一的值，却占用了完整的 32 位内存。
可以通过从 \co{->ht_resize_key} 字段
窃取一个位来消除它。
这是可行的，因为 \co{->ht_resize_key} 字段
大到足以寻址内存中的每个字节，
而 \co{ht_bucket} 结构超过一个字节长，
所以 \co{->ht_resize_key} 字段必定有
几个多余的位可用。

这种位打包技巧经常用于高度复制的数据结构中，
例如 Linux 内核中的 \co{page} 结构。
然而，可调整大小的 hash table 的 \co{ht} 结构
并没有那么高度复制。
我们应该关注的反而是 \co{ht_bucket} 结构。
缩小 \co{ht_bucket} 结构有两个主要机会：
\begin{enumerate*}[(1)]
\item 将 \tco{->htb_lock} 字段放入
\tco{->htb_head} 指针之一的低位中，以及
\item 减少所需的指针数量。
\end{enumerate*}

第一个机会可以利用 Linux 内核中的位自旋锁，
它由 \path{include/linux/bit_spinlock.h}
头文件提供。
这些在 Linux 内核中空间关键的数据结构中使用，
但并非没有缺点：

\begin{enumerate}
\item	它们比传统的自旋锁原语慢得多。
\item	它们无法参与 Linux 内核中的
	lockdep 死锁（deadlock）检测工具~\cite{JonathanCorbet2006lockdep}。
\item	它们不记录锁的所有权，
	进一步增加了调试的复杂性。
\item	它们不参与 \rt\ 内核中的优先级提升，
	这意味着持有位自旋锁时必须禁用抢占，
	这可能降低实时延迟。
\end{enumerate}

尽管有这些缺点，
当内存紧张时，位自旋锁极其有用。

第二个机会的一个方面在
\cref{sec:datastruct:Other Resizable Hash Tables}
（\cpageref{sec:datastruct:Other Resizable Hash Tables}）
中已有介绍，该节展示了只需要一组桶列表指针
的可调整大小的 hash table，
取代了
\cref{sec:datastruct:Non-Partitionable Data Structures}
（\cpageref{sec:datastruct:Non-Partitionable Data Structures}）
中介绍的可调整大小 hash table 所需的两组。
另一种方法是使用单链桶列表
来取代本章使用的双链表。
这种方法的一个缺点是删除需要额外开销，
要么通过标记待删除的指针以便稍后删除，
要么通过搜索桶列表找到要删除的元素。

简而言之，在最小内存开销与性能和简洁性之间
存在权衡。
幸运的是，现代系统上相对较大的内存
允许我们优先考虑性能和简洁性
而非内存开销。
然而，尽管 2022 年的袖珍型智能手机配备了
数千兆字节的内存，中端服务器配备了太字节级内存，
有时仍然需要采取极端措施来减少内存开销。

\subsection{硬件考虑}
\label{sec:together:Hardware Considerations}

现代计算机通常以固定大小的块
在 CPU 和主内存之间移动数据，
块大小从 32 字节到 256 字节不等。
这些块被称为\emph{\IXpl{cache line}}，
对于高性能和可扩展性极其重要，
如 \cref{sec:cpu:Overheads} 中所讨论的。
一种久经考验的扼杀性能和可扩展性的方式是
将不兼容的变量放入同一缓存行。
例如，假设一个可调整大小的 hash table 数据元素
将 \co{ht_elem} 结构与一个频繁递增的计数器
放在同一缓存行中。
频繁的递增将导致缓存行存在于
执行递增的 CPU 上，而不存在于其他任何地方。
如果其他 CPU 尝试遍历包含该元素的哈希桶列表，
它们将遭受昂贵的缓存未命中，
降低性能和可扩展性。

\begin{listing}
\begin{VerbatimL}
struct hash_elem {
	struct ht_elem e;
	long __attribute__ ((aligned(64))) counter;
};
\end{VerbatimL}
\caption{64 字节缓存行的对齐}
\label{lst:together:Alignment for 64-Byte Cache Lines}
\end{listing}

在具有 64 字节缓存行的系统上，
解决此问题的一种方法如
\cref{lst:together:Alignment for 64-Byte Cache Lines} 所示。
这里使用 \GCC 的 \co{aligned} 属性
强制 \co{->counter} 和 \co{ht_elem} 结构
位于不同的缓存行中。
这将允许 CPU 以全速遍历哈希桶列表，
尽管有频繁的递增。

当然，这引出了一个问题：
``我们怎么知道缓存行是 64 字节的？''
在 Linux 系统上，此信息可以从
\path{/sys/devices/system/cpu/cpu*/cache/} 目录获得，
甚至可以让安装过程重建应用程序
以适应系统的硬件结构。
但是，如果你希望应用程序也能在各种环境中运行，
包括非 Linux 系统，这将更加困难。
此外，即使你满足于只在 Linux 上运行，
这种自修改安装也会带来验证方面的挑战。
例如，具有 32 字节缓存行的系统可能运行良好，
但在具有 64 字节缓存行的系统上，
由于\IX{false sharing}，性能可能会受到影响。

幸运的是，有一些经验法则在实践中效果相当好，
它们被收集在 1995 年的一篇论文中~\cite{BenjaminGamsa95a}。\footnote{
	经 Orran Krieger 许可，
	这里对其中一些规则进行了转述和扩展。}
第一组规则涉及重新排列结构以适应缓存布局：

\begin{enumerate}
\item	将只读数据放在远离频繁更新数据的位置。
	例如，将只读数据放在结构的开头，
	频繁更新的数据放在结尾。
	将很少访问的数据放在中间。
\item	如果结构具有多组字段，
	每组由独立的代码路径更新，
	则将这些组彼此分开。
	同样，在组之间放置很少访问的数据可能有帮助。
	在某些情况下，将每个这样的组
	放入由原始结构引用的单独结构中
	也可能是有意义的。
\item	在可能的情况下，将更新密集型数据与
	CPU、线程或任务关联。
	我们在 \cref{chp:Counting} 的
	计数器实现中看到了这条经验法则的
	几个非常有效的例子。
\item	更进一步，按每 CPU、每线程或每任务
	的方式划分数据，如
	\cref{chp:Data Ownership} 中所讨论的。
\end{enumerate}

已有一些关于基于追踪的自动化结构字段重排的
工作~\cite{Golovanevsky:2010:TDL:2174824.2174835}。
这项工作可能会减轻从多线程软件中获得
出色性能和可扩展性所需的
最费力的任务之一。

还有一组关于锁的经验法则：

\begin{enumerate}
\item	对于保护频繁修改数据的高竞争锁，
	采取以下方法之一：
	\begin{enumerate}
	\item	将锁放在与其保护数据不同的缓存行中。
	\item	使用适合高竞争的锁，
		例如排队锁。
	\item	重新设计以减少锁竞争。
		（这种方法最好，但并不总是简单的。）
	\end{enumerate}
\item	将无竞争的锁放入与其保护数据
	相同的缓存行中。
	这种方法意味着将锁带到当前 CPU 的缓存未命中
	同时也带来了其数据。
\item	使用\IXpl{hazard pointer}、RCU 或
	对于长时间临界区使用读写锁来保护只读数据。
\end{enumerate}

当然，这些是经验法则而非绝对规则。
需要一些实验来确定哪些最适用于给定的情况。
